\documentclass{article}

\usepackage[utf8]{inputenc}
\usepackage[T1]{fontenc}
\usepackage[french]{babel}
\usepackage{amsthm}
\title{Mathématique discrètes}
\author{Loic Papon}
\date{8/10/2026}
\newtheorem{theorem}{théorème}


\begin{document}

\maketitle
\newpage
\tableofcontents
\newpage

\begin{theorem}
Pour tout entier naturel n, la somme des n premiers entiers
naturels est égal à la moitié du produit de n par (n - 1).
\end{theorem}

Soit $ n\in\mathbb { N } $ .
Notons $ S_n $ la somme des $ n $ premiers entiers naturels .
Alors :
S_n & = \sum_ { i = 0 } ^ { n - 1 } i
= 0 + 1 + \cdots + ( n - 2) + ( n - 1)
& \text { par d é finition } \\
& = \sum_ { i = 0 } ^ { n - 1 } ( n - 1 - i )
= ( n - 1) + ( n - 2) + \cdots + 1 + 0
& \text { par commutativit é } \\
En sommant les deux formes de la somme $ S_n $ donn é es ci - dessus,
nous pouvons obtenir :
2 \times S_n
& = \left ( \sum_ { i = 0 } ^ { n - 1 } i\right )
+ \left ( \sum_ { i = 0 } ^ { n - 1 } ( n - 1 - i ) \right ) \\
& = \sum_ { i = 0 } ^ { n - 1 } ( i + ( n - 1 - i ))
& \text { par associativit é et commutativit é } \\
& = \text { \emph { \ldots < compl é tez la preuve ici > \ldots } }
\end { align * }
\emph { \ldots < concluez la preuve ici > \ldots }

\end{document}